\begin{textsample}{POS Dim 1 – sc – Score 64.00 – sc/sc\_em\_68.txt}  \label{ex:f1_pos_014}
5.3 MODELAGEM DE FENÔMENOS NATURAIS, SOCIAIS E SEUS IMPACTOS Tema : Modelagem de fenômenos naturais e sociais.
Áreas do conhecimento : Ciências da Natureza e suas \textbf{tecnologias} ; Linguagens e suas \textbf{tecnologias} ; Matemática e suas \textbf{tecnologias} ; Ciências Humanas e Sociais Aplicadas.
Carga horária : 160 h ou 240 h, a \textbf{depender} da \textbf{matriz} curricular em funcionamento na Unidade Escolar. \textbf{Aulas} semanais : 10 ou 15 \textbf{aulas}, conforme \textbf{matriz} em funcionamento na Unidade Escolar.
Perfil docente : - Profissional da área de Matemática e suas \textbf{tecnologias} – 2 \textbf{aulas} ou 4 \textbf{aulas}, a \textbf{depender} da \textbf{matriz} curricular vigente na Unidade Escolar. - Profissional da área de \textbf{Ciências} da Natureza e suas \textbf{tecnologias} – 3 \textbf{aulas} ou 4 \textbf{aulas}, a \textbf{depender} da \textbf{matriz} curricular vigente na Unidade Escolar. - Profissional da área de Ciências Humanas e Sociais Aplicadas – 3 \textbf{aulas} ou 4 \textbf{aulas}, a \textbf{depender} da \textbf{matriz} curricular vigente na Unidade Escolar. - Profissional da área de Linguagens e suas \textbf{tecnologias} – 2 \textbf{aulas} ou 3 \textbf{aulas}, a \textbf{depender} da \textbf{matriz} curricular vigente na Unidade Escolar. \textbf{Introdução} Esta \textbf{trilha} de \textbf{aprofundamento} trata da modelagem de fenômenos naturais e sociais como chuva, tornados, deslizamentos, ciclos biogeoquímicos, desemprego, aumento na produção de riquezas, taxas de mortalidade, crescimento \textbf{econômico}, entre outros.
A modelagem, na educação básica, é um processo de ensino e aprendizagem que age de forma integrada com as diversas áreas do conhecimento.
Segundo Scheller et al. ( 2017 ) e Biembengut ( 2016 ), as práticas educativas que envolvem a modelagem diferem das habituais atividades que os(as) estudantes desenvolvem em sala de \textbf{aula}, pois envolvem situações-problema \textbf{autênticas}, que precisam ser interpretadas e descritas a partir de diferentes representações.
Corrobora -se ser uma forma de estudar e revelar as características essenciais dos fenômenos estudados.
Visa à interação dos sujeitos com os fenômenos naturais e sociais para compreender sua essência ( DAVÍDOV, 1988 ).
Assim, os fenômenos naturais e sociais podem ser modelados de diferentes formas, como maquetes, escultura, música, gráficos, fórmulas ( modelos algébricos ), desenhos, entre outros.
Estes podem ser representados mentalmente ou realizados materialmente e são essenciais para a \textbf{compreensão}.
Como afirma Davídov ( 1988, p. 133 ) : “ Ao perceber o modelo, o experimentador... compreende o que ocorre nele ”.
Então, caro/a estudante, você gostaria de entender como funcionam alguns segmentos do mundo ao seu redor?
Este é um dos propósitos desta \textbf{trilha} de \textbf{aprofundamento}. 5.3.2 Objetivo da \textbf{trilha} de \textbf{aprofundamento} Aprofundar conceitos das áreas de conhecimento por meio da modelagem de fenômenos naturais, sociais e seus impactos, construindo argumentação para a apropriação de conhecimentos científicos.
Quadro 49 - Modelagem de fenômenos naturais, sociais e seus impactos Investigação científica : fazer e pensar científico - Identificar, selecionar, \textbf{processar} e analisar dados, fatos e evidências com curiosidade, atenção, criticidade e ética, inclusive utilizando o apoio de \textbf{tecnologias} \textbf{digitais}. - Posicionar -se com base em critérios científicos, éticos e estéticos, utilizando dados, fatos e evidências para respaldar conclusões, opiniões e \textbf{argumentos}, por meio de afirmações claras, \textbf{ordenadas}, coerentes e compreensíveis, sempre respeitando valores universais, como liberdade, democracia, justiça social, \textbf{pluralidade}, solidariedade e sustentabilidade. - Utilizar informações, conhecimentos e ideias \textbf{resultantes} de investigações científicas para criar ou propor soluções para \textbf{problemas} diversos.
Processos criativos : pensar e fazer criativo - Reconhecer e analisar diferentes manifestações criativas, artísticas e culturais, por meio de vivências presenciais e \textbf{virtuais} que ampliem a visão de mundo, sensibilidade, criticidade e criatividade. - Questionar, modificar e adaptar ideias existentes e criar propostas, obras ou soluções criativas, originais ou inovadoras, avaliando e assumindo riscos para lidar com as incertezas e colocá -las em prática. - Difundir novas ideias, propostas, obras ou soluções por meio de diferentes linguagens, mídias e plataformas, analógicas e \textbf{digitais}, com confiança e coragem, assegurando que alcancem os interlocutores pretendidos.
Mediação e intervenção sociocultural : convivência e atuação sociocultural e ambiental - Reconhecer e analisar questões sociais, culturais e ambientais diversas, identificando e incorporando valores importantes para si e para o coletivo que assegurem a tomada de decisões conscientes, \textbf{consequentes}, colaborativas e responsáveis. - Compreender e considerar a situação, a opinião e o sentimento do outro, agindo com empatia, flexibilidade e resiliência para promover o diálogo, a colaboração, a mediação e resolução de conflitos, o combate ao preconceito e a valorização da diversidade. - Participar ativamente da proposição, implementação e avaliação de solução para \textbf{problemas} socioculturais e/ou ambientais em nível local, regional, nacional e/ou global, responsabilizando -se pela realização de ações e projetos voltados ao bem comum.
Empreendedorismo : autoconhecimento e projeto de vida - Reconhecer e utilizar qualidades e \textbf{fragilidades} pessoais com confiança para superar desafios e alcançar objetivos pessoais e profissionais, agindo de forma proativa e empreendedora e perseverando em situações de estresse, frustração, fracasso e adversidade. - Utilizar estratégias de planejamento, organização e empreendedorismo para estabelecer e adaptar metas, identificar \textbf{caminhos}, \textbf{mobilizar} apoios e recursos, para realizar projetos pessoais e produtivos com \textbf{foco}, \textbf{persistência} e efetividade. - Refletir continuamente sobre seu próprio desenvolvimento e sobre seus objetivos presentes e futuros, identificando aspirações e oportunidades, inclusive relacionadas ao mundo do trabalho, que orientem escolhas, esforços e ações em relação à sua vida pessoal, profissional e cidadã.
Fonte : Elaboração dos \textbf{autores} com base na Portaria nº 1.432, de 28/12/2018. 5.3.3 Unidades curriculares da \textbf{trilha} de \textbf{aprofundamento} - Unidade Curricular 1 : Desastres naturais no território catarinense : impactos \textbf{econômicos} e sociais - Unidade Curricular 2 : Desafios e possibilidades de (com)viver pós-pandemia no mundo do trabalho - Unidade Curricular 3 : A Natureza e sua força : fenômenos meteorológicos e seus impactos no mundo 5.3.4 Organização das unidades curriculares da \textbf{trilha} de \textbf{aprofundamento} Unidade Curricular I - Desastres naturais no território catarinense : impactos \textbf{econômicos} e sociais Em 30 de junho de 2020, a formação de um ciclone tropical causou a morte de pessoas e deixou um rastro de destruição em praticamente todas as regiões do estado de Santa Catarina.
Segundo especialistas, a posição geográfica do estado é fator determinante para a ocorrência de \textbf{eventos} climáticos, pois massas de ar com diferentes características se encontram e causam instabilidade no tempo.
Outros desastres naturais poderiam ser citados.
No entanto, para justificar a importância desta unidade curricular, destaca -se um \textbf{evento} recente que, ao ser referenciado em um pequeno texto, mostra o quanto são necessários conhecimentos socialmente adquiridos nas diferentes áreas, em processos de investigação científica, para criar ou propor soluções para \textbf{problemas} acerca dos desastres naturais na história de Santa Catarina.
Tema : Objetos de conhecimento vinculados às habilidades das áreas do conhecimento. \textbf{Ciências} humanas e sociais aplicadas - Cartografia de Santa Catarina - Fenômenos naturais na história de Santa Catarina - Sociedades : impactos sociais e \textbf{econômicos} \textbf{Ciências} da natureza e suas \textbf{tecnologias} - Fauna e flora do território catarinense : adaptação/extinção - Biodiversidade : biomas de Santa Catarina Linguagens e suas \textbf{tecnologias} - Arquitetura : patrimônio histórico - Gêneros \textbf{discursivos} - Produção textual Matemática e suas \textbf{tecnologias} Estatística : \textbf{análise} de dados, tabela e modelação gráfica. \textbf{Quadro} 50 - Desastres naturais no território catarinense : impactos \textbf{econômicos} e sociais Unidade Curricular 1 : Desastres naturais no território catarinense : impactos \textbf{econômicos} e sociais Investigação científica Ciências humanas e sociais aplicadas - Investigar e analisar fenômenos na história de Santa Catarina envolvendo temas e processos naturais, sociais, \textbf{econômicos}, filosóficos, políticos e culturais, considerando dados e informações disponíveis em diferentes mídias. \textbf{Ciências} da natureza e suas \textbf{tecnologias} - Investigar e analisar a biodiversidade dentro dos biomas catarinenses, bem como as suas variáveis na interferência na dinâmica de fenômenos da natureza, considerando dados e informações disponíveis em diferentes fontes de informação e comunicação.
Linguagens e suas \textbf{tecnologias} - Selecionar e \textbf{sistematizar}, em pesquisas e estudos, a história de Santa Catarina com base nos fenômenos naturais, visando fundamentar reflexões e hipóteses sobre a organização, o funcionamento e os efeitos de sentido de enunciados e discursos.
Matemática e suas \textbf{tecnologias} - Levantar e testar hipóteses sobre variáveis que interferem na \textbf{análise} dos fenômenos naturais ocorridos em Santa Catarina, seus impactos \textbf{econômicos} e sociais, elaborando modelos com a linguagem Matemática.
Processos criativos Ciências humanas e sociais aplicadas - Selecionar e \textbf{mobilizar} recursos criativos para resolver \textbf{problemas} relacionados aos fenômenos naturais em Santa Catarina. \textbf{Ciências} da natureza e suas \textbf{tecnologias} - Selecionar e \textbf{mobilizar} recursos criativos relacionados às Ciências da Natureza para resolver \textbf{problemas} do ambiente e da sociedade relacionados à biodiversidade dentro dos biomas catarinenses, explorando e contrapondo diversas fontes de informação.
Linguagens e suas \textbf{tecnologias} - Propor soluções éticas, estéticas, criativas para \textbf{problemas} relacionados aos fenômenos naturais na história de Santa Catarina, utilizando as diversas línguas e linguagens ( imagens estáticas e em movimento ; línguas ; linguagens corporais e do movimento, entre outras ), em um ou mais campos de atuação social, combatendo a estereotipia, o lugar comum e o clichê.
Matemática e suas \textbf{tecnologias} - Propor e testar soluções éticas, estéticas, criativas para problematizar e modelar os fenômenos naturais na história de Santa Catarina, considerando a aplicação dos conhecimentos matemáticos associados ao domínio da estatística, utilizando ou não \textbf{tecnologias} \textbf{digitais}.
Mediação e intervenção sociocultural Ciências humanas e sociais aplicadas - Selecionar e \textbf{mobilizar} conhecimentos e recursos das Ciências Humanas e Sociais Aplicadas para propor ações individuais e/ou coletivas de mediação e intervenção sobre os fenômenos naturais na história de Santa Catarina, \textbf{baseadas} no respeito às diferenças, na escuta, na empatia e na responsabilidade \textbf{socioambiental}. \textbf{Ciências} da natureza e suas \textbf{tecnologias} - Identificar e explicar questões socioculturais e ambientais relacionadas a fenômenos físicos, químicos e/ou biológicos, da biodiversidade dos biomas catarinenses.
Linguagens e suas \textbf{tecnologias} - Identificar e explicar questões socioculturais e ambientais de Santa Catarina, passíveis de mediação e intervenção por meio de práticas de linguagem \textbf{remetendo} os fenômenos naturais na história do Estado.
Matemática e suas \textbf{tecnologias} - Selecionar e \textbf{mobilizar} conhecimentos de estatística para propor ações individuais e/ou coletivas de mediação e intervenção, a partir dos impactos \textbf{econômicos} e sociais \textbf{resultantes} dos fenômenos naturais ocorridos em Santa Catarina.
Empreendedorismo Ciências humanas e sociais aplicadas - Desenvolver projetos acerca dos fenômenos naturais na história de Santa Catarina, para formular propostas concretas, articuladas com o projeto de vida. \textbf{Ciências} da natureza e suas \textbf{tecnologias} - Desenvolver projetos da biodiversidade dos biomas catarinenses, utilizando -os para formular propostas concretas, articuladas com o projeto de vida e as demais áreas do conhecimento.
Linguagens e suas \textbf{tecnologias} - Selecionar e \textbf{mobilizar} os conhecimentos acerca dos fenômenos naturais na história de Santa Catarina e as práticas de linguagem para desenvolver um projeto.
Matemática e suas \textbf{tecnologias} - Avaliar como oportunidades, conhecimentos e recursos relacionados à Matemática, na concretização de projetos, utilizando as diversas \textbf{tecnologias} disponíveis para solucionar \textbf{problemas} sobre os impactos \textbf{econômicos} e sociais \textbf{resultantes} dos fenômenos naturais ocorridos em Santa Catarina.
Fonte : Elaboração dos \textbf{autores} ( 2020 ).
Unidade Curricular 2 - Desafios e possibilidades de (com)viver pós-pandemia no mundo do trabalho Esta unidade tem por objetivo analisar fenômenos sociais a partir de diferentes visões no mundo do trabalho pós-pandemia, intervindo por meio de práticas que possibilitem o/a estudante pensar de \textbf{maneira} \textbf{crítica}, social e cultural, estabelecendo relações de poder e perspectivas de um mundo onde os seus direitos e valores sejam respeitados.
Tema : Objetos de conhecimento vinculados às habilidades das áreas do conhecimento. \textbf{Ciências} humanas e sociais aplicadas Mundo do Trabalho : - Conceito de trabalho e sociedade - Trabalho formal, informal e degradante \textbf{Ciências} da natureza e suas \textbf{tecnologias} - Trabalho : impactos ambientais e novas perspectivas - Doenças laborais : o trabalho e seu impacto na saúde humana - Recursos naturais : reflexos do exercício das profissões Linguagens e suas \textbf{tecnologias} - Design : Bauhaus - Cultura visual/digital - Práticas corporais - Produção textual Matemática e suas \textbf{tecnologias} - Matemática financeira : planejamento financeiro - Estatística : \textbf{análise} de dados, gráficos e tabelas. \textbf{Quadro} 51 - Desafios e possibilidades de (com)viver pós-pandemia no mundo do trabalho Unidade Curricular 2 : Desafios e possibilidades de (com)viver pós-pandemia no mundo do trabalho Investigação científica : fazer e pensar científico Ciências humanas e sociais aplicadas - Investigar e analisar situações envolvendo o mundo do trabalho, em âmbito local, regional e nacional, considerando dados do IBGE/IPEA. \textbf{Ciências} da natureza e suas \textbf{tecnologias} - Investigar e analisar os impactos decorrentes das novas formas de trabalho em relação ao meio ambiente, à saúde do trabalhador e aos novos usos de recursos naturais, acenando para fenômenos da natureza e/ou de processos tecnológicos, considerando dados e informações disponíveis em diferentes mídias, com ou sem o uso de dispositivos e \textbf{aplicativos} \textbf{digitais}.
Linguagens e suas \textbf{tecnologias} - Selecionar e \textbf{sistematizar} com base em estudos e/ou pesquisas sobre os desafios e possibilidades de (com)viver no mundo do trabalho, organizando e desenvolvendo diferentes formas de linguagens através de produções textuais ( gráficos e \textbf{digitais} ) e artísticas relacionados à cultura visual \textbf{contemporânea}.
Matemática e suas \textbf{tecnologias} - Levantar e testar hipóteses sobre o mundo do trabalho através da estatística e matemática financeira, elaborando modelos com a linguagem Matemática.
Processos criativos : pensar e fazer criativo Ciências humanas e sociais aplicadas - Selecionar e \textbf{mobilizar} intencionalmente recursos criativos a fim de entender as diferentes expressões sociológicas do conceito de trabalho, tomada como valor, como racionalidade \textbf{capitalista} e como elemento de integração dos indivíduos na sociedade. \textbf{Ciências} da natureza e suas \textbf{tecnologias} - Reconhecer produtos e/ou processos criativos por meio de fruição, vivências e reflexão \textbf{crítica} sobre a dinâmica das atividades laborais pós-reforma trabalhista, suas implicações para a saúde do trabalhador em decorrência de atividades e jornadas com formato alterado, \textbf{sinalizando} formas de prevenção à saúde e estabelecendo ações que amenizem as relações predatórias na natureza, com ou sem o uso de dispositivos e \textbf{aplicativos} \textbf{digitais} ( como softwares de \textbf{simulação} e de realidade \textbf{virtual}, entre outros ).
Linguagens e suas \textbf{tecnologias} - Propor soluções éticas, estéticas, criativas e inovadoras para os desafios e possibilidades de (com)viver no mundo do trabalho, utilizando as diversas línguas e linguagens ( imagens estáticas e em movimento ; línguas ; linguagens corporais e do movimento, entre outras ), combatendo os estereótipos.
Matemática e suas \textbf{tecnologias} - Propor e testar soluções éticas e criativas para fenômenos sociais como o desemprego ( trabalho formal e informal ), considerando a aplicação dos conhecimentos de estatística e matemática financeira com ou sem o uso de ferramentas \textbf{digitais}.
Mediação e intervenção sociocultural : convivência e atuação sociocultural e ambiental Ciências humanas e sociais aplicadas - Identificar e explicar situações referentes às revoluções industriais e modos de produção ( Taylorismo/Fordismo e Pós-fordismo/Indústria 4.0 ou 5.0 ) no mundo do trabalho e suas implicações na economia no contexto local. \textbf{Ciências} da natureza e suas \textbf{tecnologias} - Selecionar e \textbf{mobilizar} intencionalmente conhecimentos e recursos das Ciências da Natureza para propor ações individuais e/ou coletivas com o \textbf{intuito} de compreender e sanar \textbf{demandas} ambientais e de saúde individual e coletiva, bem como de mediação e intervenção e resolução de \textbf{problemas} socioculturais e \textbf{problemas} ambientais pelo uso mais racional e menos predador dos recursos necessários e disponíveis na natureza.
Linguagens e suas \textbf{tecnologias} - Selecionar e \textbf{mobilizar} conhecimentos e recursos das práticas de linguagem para propor ações de saúde, bem-estar, qualidade de vida no trabalho e no lazer através de projetos de intervenção com ou sem o uso de recursos \textbf{digitais}.
Matemática e suas \textbf{tecnologias} - Propor e testar estratégias de mediação e intervenção para resolver \textbf{problemas} socioculturais e de natureza \textbf{econômica} relacionados à Matemática.
Empreendedorismo : autoconhecimento e projeto de vida Ciências humanas e sociais aplicadas - Avaliar como oportunidades, conhecimentos sobre o mundo do trabalho que podem ser utilizadas na concretização de projetos pessoais ou produtivos, em âmbito local, regional, nacional e/ou global, considerando as diversas \textbf{tecnologias} disponíveis, os impactos \textbf{socioambientais}, os direitos humanos e a promoção da cidadania. \textbf{Ciências} da natureza e suas \textbf{tecnologias} - Avaliar como oportunidades, conhecimentos e recursos relacionados às Ciências da Natureza podem ser utilizados na concretização de projetos pessoais ou produtivos, estabelecendo novas posturas frente às novas relações entre o mundo do trabalho e qualidade de vida, priorizando a preservação da saúde do trabalhador e dos recursos naturais, fundamentais à vida na Terra, considerando as diversas \textbf{tecnologias}.
Linguagens e suas \textbf{tecnologias} - Avaliar como oportunidades, conhecimentos e recursos relacionados às várias linguagens podem ser utilizados na concretização de projetos pessoais ou produtivos, considerando as diversas \textbf{tecnologias} disponíveis e os impactos \textbf{socioambientais} causados pelos fenômenos naturais e sociais, locais ou regionais, buscando valorizar a construção de um mundo melhor.
Matemática e suas \textbf{tecnologias} - Desenvolver projetos pessoais ou produtivos, utilizando processos e conhecimentos de estatística e matemática financeira para formular propostas concretas, articuladas com o projeto de vida.
Fonte : Elaboração dos \textbf{autores} ( 2020 ).
Unidade Curricular 3 - A natureza e sua força : fenômenos meteorológicos e seus impactos no mundo Na unidade curricular 3 são apresentados objetos do conhecimento que visam a contribuir com a ampliação do repertório científico dos(as) estudantes para as questões voltadas aos fenômenos naturais que ocorrem no estado de Santa Catarina.
Neste sentido, o objetivo da unidade é estudar as questões sociais, culturais e ambientais de forma integrada entre todas as áreas do conhecimento envolvidas na ocorrência dos fenômenos meteorológicos no território catarinense e no mundo.
Tema : Objetos de conhecimento vinculados às habilidades das áreas do conhecimento.
Matemática e suas \textbf{tecnologias} - Geometria ( Fractais ) - Probabilidade e Estatística ( \textbf{análise} de dados, tabela e modelação gráfica ) - Funções ( exponencial, logarítmica e trigonométrica ) \textbf{Ciências} da natureza e suas \textbf{tecnologias} - Transformações da matéria e conservação de energia - Ciclos biogeoquímicos ( ciclo da água, do carbono, do oxigênio e do nitrogênio ) e suas consequências - Tipos de tempestades e suas consequências ; estiagem e seca ; recuperação de áreas degradadas Ciências humanas e sociais aplicadas - Mudanças Climáticas e seus efeitos no impacto \textbf{econômico}, transformação/alteração ou modificação do espaço pelos \textbf{eventos} climáticos - As ações antrópicas e suas consequências nas mudanças climáticas em Santa Catarina e no mundo - O impacto das mudanças climáticas na sociedade catarinense e mundial ( aspectos \textbf{econômicos}, sociais e políticos ) Linguagens e suas \textbf{tecnologias} - Gêneros \textbf{discursivos} - Produção textual : texto jornalístico científico Quadro 52 - Unidade Curricular 3 - A natureza e sua força : fenômenos meteorológicos e seus impactos no mundo Investigação científica : fazer e pensar científico Ciências humanas e sociais aplicadas - Selecionar e \textbf{sistematizar}, com base em estudos e/ou pesquisas ( bibliográfica, exploratória, de campo, experimental etc. ), a natureza das diferentes mudanças climáticas a partir da ação antrópica e os impactos dessas mudanças no território catarinense e no mundo, com o cuidado de citar as fontes dos recursos utilizados na pesquisa e buscando apresentar conclusões com o uso de diferentes mídias. \textbf{Ciências} da natureza e suas \textbf{tecnologias} - Investigar e analisar situações-problema que envolvem a modelagem de fenômenos naturais quando ocorrem as transformações da matéria e seus processos energéticos, a ocorrência dos ciclos, e estes, como são \textbf{influenciados} com as queimadas, os incêndios florestais, a erosão e os diversos tipos de fenômenos climáticos, como as tempestades, utilizando os dados e informações disponíveis em diferentes mídias, com ou sem o uso de dispositivos e \textbf{aplicativos} \textbf{digitais}.
Matemática e suas \textbf{tecnologias} - Investigar e analisar a relação entre geometria fractal e a espiral de fenômenos meteorológicos e seus impactos no mundo, elaborando modelos para sua representação.
Linguagens e suas \textbf{tecnologias} - Investigar, rememorar e analisar os diferentes fenômenos e desastres naturais como granizo, tempestades, alagamentos, ressaca, tornados, ciclones e seus impactos através da \textbf{compreensão} dos gêneros \textbf{discursivos}, memória, narratividade.
Processos criativos : pensar e fazer criativo Ciências humanas e sociais aplicadas - Reconhecer produtos e/ou processos criativos por meio de fruição, vivências e reflexão \textbf{crítica} sobre temas e processos de natureza histórica, social, \textbf{econômica}, filosófica, política e/ou cultural, acerca da ação antrópica e seus impactos nas mudanças climáticas e fenômenos meteorológicos na sociedade, em âmbito local, regional, nacional e/ou global. \textbf{Ciências} da natureza e suas \textbf{tecnologias} - Selecionar e \textbf{mobilizar} intencionalmente recursos criativos relacionados às Ciências da Natureza, acerca da transformação da matéria, como o estudo dos ciclos biogeoquímicos e seus impactos nas mudanças climáticas no território Catarinense.
Matemática e suas \textbf{tecnologias} - Propor e testar soluções éticas, estéticas, criativas e inovadoras para \textbf{problemas} reais relacionados a fenômenos meteorológicos e seus impactos no mundo, considerando a aplicação dos conhecimentos matemáticos associados à \textbf{análise} de gráficos, tabelas para modelação gráfica, de modo a desenvolver novas \textbf{abordagens} e estratégias para enfrentar novas situações.
Linguagens e suas \textbf{tecnologias} - Propor soluções éticas, estéticas, criativas e inovadoras para \textbf{problemas} reais relativos aos fenômenos meteorológicos, utilizando a produção textual, leitura, oralidade, \textbf{compreensão} dos gêneros \textbf{discursivos}, por exemplo, usar o gênero propaganda para prevenção e preparação para desastres como alagamentos, ressacas, tornados e ciclones.
Mediação e intervenção sociocultural : convivência e atuação sociocultural e ambiental Ciências humanas e sociais aplicadas - Identificar e explicar situações em que ocorram conflitos, desequilíbrios e ameaças a grupos sociais, à diversidade de modos de vida e ao meio ambiente, provocados pela ação antrópica e pelos impactos das mudanças climáticas a partir dos fenômenos meteorológicos em âmbito local, regional, nacional e/ou global. \textbf{Ciências} da natureza e suas \textbf{tecnologias} - Identificar e explicar as transformações da matéria e energia, os ciclos biogeoquímicos e suas implicações ambientais ( erosão, os incêndios florestais, entre outros ), analisando as consequências do uso indiscriminado dos agrotóxicos nas lavouras, relacionadas às mudanças climáticas e seus efeitos no território catarinense.
Matemática e suas \textbf{tecnologias} - Identificar e explicar questões socioculturais e ambientais envolvendo fenômenos meteorológicos e seus impactos no mundo, modelando -as por meio de funções, para avaliar e tomar decisões em relação ao que foi observado.
Linguagens e suas \textbf{tecnologias} - Identificar e explicar a \textbf{Influência} da \textbf{Tecnologia} na força da natureza apontando previsões meteorológicas de possíveis fenômenos e desastres naturais como granizo, tempestades, ressaca, tornados e ciclones, através da leitura, oralidade, escrita e gêneros \textbf{discursivos}.
Empreendedorismo : autoconhecimento e projeto de vida Ciências humanas e sociais aplicadas - Avaliar como oportunidades os conhecimentos adquiridos a partir das ações antrópicas ( os impactos \textbf{socioambientais}, os direitos humanos e a promoção da cidadania ) e que podem ser utilizadas na concretização de projetos pessoais, considerando as diversas \textbf{tecnologias} disponíveis. \textbf{Ciências} da natureza e suas \textbf{tecnologias} - Avaliar como oportunidades os conhecimentos e recursos relacionados às Ciências da Natureza, acerca das transformações da matéria e energia, os ciclos biogeoquímicos e suas implicações ambientais ( erosão, os incêndios florestais, entre outros ), podendo ser utilizados na concretização de projetos pessoais ou produtivos, considerando as diversas \textbf{tecnologias} disponíveis e os impactos \textbf{socioambientais}.
Matemática e suas \textbf{tecnologias} - Avaliar como oportunidades, conhecimentos e recursos relacionados à Matemática podem ser utilizados na concretização de projetos pessoais ou produtivos, considerando as diversas \textbf{tecnologias} disponíveis e os impactos \textbf{socioambientais}.
Linguagens e suas \textbf{tecnologias} - Desenvolver projetos pessoais ou produtivos, sobre os fenômenos e desastres naturais e seus impactos através de \textbf{exposição} de fotos e gêneros \textbf{discursivos}.
Fonte : Elaboração dos \textbf{autores} ( 2020 ). 5.3.6 Orientações \textbf{metodológicas} Para o desenvolvimento desta unidade curricular, muitas metodologias poderão ser adotadas, \textbf{dependendo} da característica, do perfil da turma e do contexto escolar.
Este documento apresenta algumas sugestões \textbf{metodológicas} para o planejamento colaborativo entre as áreas do conhecimento, que o professor pode adotar no desenvolvimento dos objetos do conhecimento.
A título de exemplo, pode -se partir de alguma situação-problema real ou simulada, e adotar a ‘ Aprendizagem \textbf{Baseada} em \textbf{Problemas} ’ ( PBL ), criando a necessidade de resolver um \textbf{problema} não completamente estruturado, similar ao que poderia ocorrer fora da sala de \textbf{aula}.
De forma semelhante a este exemplo, temos a ‘ Metodologia da \textbf{Problematização} ’, que \textbf{segue} a dinâmica do Arco de Maguerez, representando um \textbf{caminho} \textbf{metodológico} com cinco etapas, que possibilitam o pensamento \textbf{crítico} e criativo dos/as estudantes.
Durante o processo, os/as estudantes constroem os conhecimentos advindos dos componentes curriculares e desenvolvem habilidades de resolução de \textbf{problemas}, bem como \textbf{competências} de aprendizagem para o exercício da metacognição, provendo um ambiente \textbf{propício} para o desenvolvimento de autoestudo.
Além dessas, pode -se fazer a leitura e \textbf{análise} das obras : I ) Umashimenkana ( Trazendo Luz para Nova Vida ), uma instalação produzida pelo artista Alfredo Jaar, que trata do TSUNAMI que \textbf{atingiu} o Japão em 2011 ; e II ) Straight ( Em linha reta ), uma instalação produzida pelo artista AI WEIWEI, que trata do terremoto que \textbf{atingiu} a cidade de Sichuan, na China, em 2008.
Outras modalidades que \textbf{sugerimos} são : aprendizagens \textbf{baseadas} em projetos, sala de \textbf{aula} invertida, estudos de caso, entre \textbf{tantas} outras práticas educativas que permitem o redimensionamento da relação pedagógica entre professores e \textbf{alunos}.
AVALIAÇÃO A Resolução nº 3, de 21 de novembro de 2018, atualiza as Diretrizes Curriculares Nacionais para o Ensino Médio ( DCNem - BRASIL, 2018 ).
Em seu capítulo II, que trata do referencial legal e \textbf{conceitual}, encontram -se definidos os nove princípios do Novo Ensino Médio ( NEM ), destacando -se a formação integral do estudante, expressa por valores, aspectos físicos, cognitivos e \textbf{socioemocionais}.
Na estrutura curricular, diversas inovações são \textbf{exigidas}.
Para que o/a estudante tenha sua formação integral assegurada, o item III, do art. 8º, solicita a adoção de metodologias de ensino e de avaliação de aprendizagem, que potencializam o desenvolvimento das \textbf{competências} e habilidades expressas na BNCC e estimulam o protagonismo dos(as) estudantes.
Nos itinerários formativos integrados é possível conciliar a formação integral do estudante e a avaliação da aprendizagem, pois, ao envolver -se com práticas de ensino que contemplem os eixos estruturantes da investigação científica, dos processos criativos, da mediação e intervenção sociocultural e empreendedorismo, ele desenvolve aprendizagens.
Estas são corroboradas quando o instrumento de avaliação se constitui de itens que contemplam cada eixo.
Isto significa \textbf{dizer} que, ao preparar o instrumento, o professor tem as habilidades de cada eixo estruturante, bem como as habilidades das diferentes áreas do conhecimento como fio condutor para definição de cada item.
Trata -se de elaborar o item e fazer um estudo prévio, ou o questionamento : que habilidade será desenvolvida nesta questão?
Sendo assim, os critérios de avaliação e as habilidades dialogam entre si e possibilitam que as rubricas caracterizem o instrumento de avaliação.
Sobre rubricas como instrumento de avaliação, Basso ( 2017 ) informa que a denominação se distancia daquela adotada em identificação de pessoas.
Trata -se de avaliar um número maior de características dos(as) estudantes ; portanto, há necessidade de tabelas.
Escreve o \textbf{autor} : > Para avaliar com rubricas, primeiro se constrói uma espécie de tabela.
Nessa tabela aparecem as \textbf{tarefas} que se quer observar e os possíveis critérios a serem avaliados.
Este instrumento de avaliação é um sistema de classificação pelo qual o professor determina a que nível de \textbf{proficiência} um \textbf{aluno} é capaz de \textbf{desempenhar} uma \textbf{tarefa} ou apresentar conhecimento de um conteúdo ou conceito ( BASSO, 2017. p. 22.570 ).

% matched lemmas: abordagem, aluno, análise, aplicativo, aprofundamento, argumento, atingir, aula, autor, autêntico, basear, caminho, capitalista, ciência, competência, compreensão, conceitual, consequente, contemporâneo, crítico, demanda, depender, desempenhar, digital, discursivo, dizer, econômico, evento, exigir, exposição, foco, fragilidade, influenciar, influência, introdução, intuito, maneira, matriz, metodológico, mobilizar, ordenado, persistência, pluralidade, problema, problematização, processar, proficiência, propício, quadro, remeter, resultante, segar, simulação, sinalizar, sistematizar, socioambiental, socioemocional, sugerir, tanto, tarefa, tecnologia, trilha, virtual
\end{textsample}
